% ==========================================================================
%  Chapter 103 -- CA Meta-Tuning of the Continuation Solver (Extension, §P3)
% ==========================================================================

\chapter{Cultural-Algorithm Meta-Tuning of the Continuation Solver}
\label{ch:ca-tuner-solver}

\paragraph{Normative status.}
Phase~P3 of the branch-selection extension
(Cap.~\ref{ch:branch-selection-reversals}). The meta-tuning \emph{mechanism}
--- the Cultural-Algorithm module (Cap.~\ref{ch:cultural-algorithm-module})
driving the continuation solver's configuration
(Cap.~\ref{ch:deflated-continuation}) --- is implemented and verified in
isolation by ctest \texttt{regularized\_newton\_ca\_tuner}
(\path{tests/test_regularized_newton_ca_tuner.cpp}). Status:
\texttt{mechanism\_verified\_in\_isolation}. Coupling the tuner to a
window replay of the real cyclic driver, and the full Gate~G4 evidence,
belong to the driver integration (Phase~I).

\section{Solver configuration as a decision vector}

The three remedies of Caps.~\ref{ch:deflated-continuation}
and~\ref{ch:energy-tao-continuation} change \emph{how} a step is taken.
Phase~P3 is orthogonal: it tunes the numbers the continuation already
exposes so that the plain LM step succeeds where a default choice fails.
Those numbers form a decision vector
\(\mathbf{g}=(\log_{10}\mu_0,\ \text{grow},\ \text{drop},\
f_{\mu_{\max}},\ n_{\text{stag}},\dots)\)
over a box of admissible values. The map from \(\mathbf{g}\) to an
outcome is cheap to evaluate (replay one reversal) but has no usable
derivative --- it threads through a nonsmooth accept/reject loop --- which
is exactly the setting the Cultural Algorithm was built for
(Cap.~\ref{ch:cultural-algorithm-module}, \cite{Reynolds1994Cultural}).

\section{Objective: a replay of the continuation}

The objective evaluates a genome by \emph{replaying} the solver with the
decoded configuration and scoring the result:
\[
  \text{fitness}(\mathbf{g})
  = -\bigl\lVert\mathbf{R}\bigl(\mathbf{u}^{\text{end}}(\mathbf{g})\bigr)\bigr\rVert,
\]
maximised, so the search drives the end-of-step residual to zero. In the
real driver the replay runs a short window \(k_0,\dots,k_0{+}m{-}1\) across
a reversal from a captured checkpoint, and the fitness combines the
converged fraction with a spike penalty
\(\max\!\bigl(0,\ |V|_{\text{peak}}/|V|_{\text{env}}-1\bigr)\); the
checkpoint/restore that brackets each candidate is the
\texttt{capture\_state}/\texttt{restore\_state} pair added in Phase~P1
(Cap.~\ref{ch:deflated-continuation}). The isolated test uses the same
shape on a 1-DOF surrogate so the tuner is validated independently of the
column's cost.

\section{An isolated proof that tuning matters}

The surrogate is the cubic \(R(u)=(u-1)(u-2)(u-4)\) started at
\(u_0=3.5\). The default LM configuration \emph{stalls} there: the Newton
step overshoots root~\(4\) to \(u\approx4.5\), where \(|R|\) is larger,
and the monotone-residual acceptance rejects it
(Cap.~\ref{ch:deflated-continuation} discusses this same overshoot). More
regularisation shrinks the step: with \(\mu_{\max}=f_{\mu_{\max}}\lVert
\operatorname{diag}\mathbf{K}\rVert\) large enough the trial lands below
\(4.5\), the residual decreases and the iteration accepts and converges.
The Cultural Algorithm searches the five-parameter box and \emph{finds}
such a configuration:
%
\begin{center}
\begin{tabular}{@{}lll@{}}
\toprule
Configuration        & end-of-step \(|R|\)          & outcome \\
\midrule
Default              & \(1.875\)                    & stalls \\
CA-tuned             & \(8.9\times10^{-12}\)        & converges \\
\bottomrule
\end{tabular}
\end{center}
%
The winning genome raises \(f_{\mu_{\max}}\) from the default \(0.1\) to
\(\approx0.69\) --- precisely the ``more regularisation'' the physics
predicts --- and the search reaches it within a single generation. The
run is bit-for-bit deterministic under a fixed seed. This closes the
mechanism: the tuner is not a black box hoping for luck, it recovers the
qualitatively-correct knob.

\section{Genome v1 for the real driver}

The production genome is eight-dimensional and \emph{solver-only}, since
the material-latched genes (closure stiffness fraction, anchor cap) are
read once into \texttt{static const} and cannot be retuned in-process:
\[
  \bigl(\log_{10}\mu_0,\ \text{grow},\ \text{drop},\ n_{\text{stag}},\
   s_{\text{pred}}^{\max},\ \log_{10}\varepsilon_{\text{accept}},\
   n_{\text{rev-subdiv}},\ n_{\text{rev-sub}}\bigr).
\]
Each candidate is applied through \texttt{set\_config}/%
\texttt{set\_regularization} (Cap.~\ref{ch:deflated-continuation}) without
rebuilding the analysis. The reversal-subdivision genes are included even
though prior evidence suggests subdivision \emph{worsens} the spurious
branch: leaving them in the search is how the tuner is asked to
\emph{confirm} that, rather than assuming it.

\section{Honest budget and scope}

A replay campaign is not free: a population of \(\sim\!8\) over
\(\sim\!10\) generations replaying a \(\sim\!6\)-step window is
\(\mathcal{O}(500)\) continuation steps, minutes on the
\(2\times2\times4\) mesh, and it must never run inside the
\(5\times5\times20\) mesh. The tuner therefore runs once, offline, on the
reduced mesh, and its winning genome is promoted to a full run and
compared against \texttt{g0\_baseline}. The material-latched genes are
swept by an outer per-process loop, not by the CA. These constraints, and
the Gate~G4 criterion (regularized reversal steps cut by half, or peak
\(\le1.5\times\) envelope, with a reproducible non-decreasing fitness
history), are the subject of the Phase~I campaign. The same
\texttt{maximize} will later tune control devices rather than solver knobs
(Cap.~\ref{ch:cultural-algorithm-module},
\cite{LaraValenciaEchavarriaValencia2024}); nothing in the tuner is
specific to the continuation problem beyond the objective it is handed.
